\section{Vertical Profile Contract}

A vertical profile defines how a class of entities expresses institutional
potential. It supplies configuration to the universal engine without changing
the evidence, provenance, or score-result contracts.

\subsection{Profile Shape}

\begin{lstlisting}
{
  "profile_id": "ncra-v1",
  "vertical_family": "consumer_facing",
  "entity_class": "nationwide_consumer_reporting_agency",
  "legal_regimes": ["FCRA", "Regulation V"],
  "peer_class": { "id": "ncra-primary", "minimum_size": 3 },
  "dimensions": {
    "F": { "metrics": ["LMCR", "PAS", "RVR", "SOF"] },
    "C": { "metrics": ["OVI", "PAR", "CDS", "DRR"] },
    "T": { "metrics": ["RDR", "RTI", "FCR", "PAX"] },
    "E": { "metrics": ["CRR", "PTS", "RUI", "BCR"] },
    "Q": { "metrics": ["DDI", "GAS", "PCOR", "ASR"] },
    "Ac": { "metrics": ["SCR", "AC", "CRR_c", "RAR"] }
  },
  "normalization_version": "ncra-norm-v1",
  "sufficiency_rules": {},
  "hard_floors": {},
  "caps": {}
}
\end{lstlisting}

A corrections profile may use different metric codes, sources, peer classes,
and thresholds while retaining the same six-dimensional result contract. The
engine must reject incomplete profiles rather than silently substituting
metrics from another vertical.

\subsection{Entity and Activity Separation}

An entity profile identifies the subject; an activity record identifies what
that subject does. This distinction is mandatory for corporate groups and
public departments with multiple programs.

\begin{lstlisting}
entity -> legal entity -> activity -> jurisdiction -> legal regime -> evidence
\end{lstlisting}

A label such as ``debt collector'' must therefore be attached only when the
source supports the named legal entity, activity, jurisdiction, and applicable
period. It must not be generalized across a corporate group without evidence.
